\chapter{The driver, the safety chain and the flight zone}
\label{ch:driver}
\setlength{\emergencystretch}{3em}

This chapter covers the node that talks to the drone, \code{tello\_driver}, and the three plain-Python modules it
is built on: \file{core/safety.py} (the chain of guards every command goes through), \file{core/zone.py} (the obstacle-free
flight zone) and \file{core/video\_reader.py} (an optional low-latency video decoder). It ends with the fake drone
(\file{sim/fake\_tello.py}, \file{sim/fake\_world.py}) that replaces the Tello when the launch file is started with
\code{mock:=true}, and the small \file{core/ros\_env.py} helper.

Every number quoted here comes from the code or from the configuration files at the time of writing. The line
numbers in the code review refer to the same version.

\section{Why there is a driver, and what goes into it}
\label{sec:driver-role}

\subsection{The single gate}

The Tello is reached over its own Wi-Fi with a text protocol on UDP (commands on port 8889, telemetry on 8890,
H.264 video on 11111). The Python library \code{djitellopy} wraps that protocol. The Tello accepts commands only from
the first program that talked to it after power-on, so exactly one process may own the link. In this project
that process is the ROS 2 node \code{tello\_driver} (\file{src/tello\_driver/tello\_driver/driver.py}).

Every other program, the keyboard teleop and all the missions (\code{tello\_goto}, \code{tello\_marker\_tour},
\code{tello\_patrol}, \code{tello\_strike}, \code{tello\_step\_test}), can only \emph{ask}: it publishes a velocity on
\topic{/tello/cmd\_vel\_safe} and calls the services \topic{/tello\_driver/takeoff}, \topic{/tello\_driver/land} and
\topic{/tello\_driver/emergency}. The driver decides what is really sent to the drone.

\begin{keypoint}
The driver is the only path from a ROS program to the motors, so whatever it enforces holds for \emph{every}
client, including a buggy mission or a human pressing the wrong key. The flight-zone check, the ceiling, the
height band, the wall check and the position hold therefore live in the driver. No mission can switch them off,
because no mission has another way to reach the drone.
\end{keypoint}

\subsection{What goes into the driver, and what goes elsewhere}

This is the rule the code follows, with examples:

\begin{center}
\begin{tabularx}{\linewidth}{>{\raggedright\arraybackslash}p{0.34\linewidth}X}
\toprule
\textbf{In the driver} & \textbf{Why} \\
\midrule
The \code{djitellopy} link, connection checks, \code{streamon}, video watchdog & only one process may own the link \\
Conversion m/s $\rightarrow$ rc joystick points, axis signs & a property of the drone, not of a mission \\
Command timeout (0.5 s) and the 20 Hz rc loop & the Tello keeps flying with the last rc value it received \\
Takeoff, the climb to the marker height, landing, low battery & must work for every client \\
Position hold, ceiling, height band, flight zone, wall ahead (\file{core/safety.py}) & must protect every client \\
Publishing \topic{/tello/rc\_sent} & the EKF needs what was \emph{really} sent, after the guards \\
\midrule
\textbf{Elsewhere} & \\
\midrule
State estimation (\file{core/ekf\_node.py}, \file{core/ekf\_core.py}) & heavy, testable on recorded flights (replay) \\
Marker detection and PnP (\file{aruco\_localizer.py}) & vision, separate package \\
Where to go: goals, paths, search behaviour (missions) & changes per task; may be wrong without risk \\
Visualisation of the zone (\file{marker\_map\_publisher.py}) & RViz only, no effect on the flight \\
\bottomrule
\end{tabularx}
\end{center}

Inside the driver there is a second split. \file{core/driver.py} contains the ROS wiring and the drone link;
\file{core/safety.py} (class \code{CommandFilter}) contains every decision about the commands and imports no ROS module.
The driver feeds it numbers (pose, velocity, drift, height, marker distance) and asks for the filtered command.
The filter can therefore be run in a plain Python loop against a simulated drone; its docstrings mention such
"stress tests" (seeds 27, 142, 203, 367), but the test script itself is not in the repository (see the code review).
\file{core/zone.py} is ROS-free for the same reason, and the missions import it too.

\subsection{Interfaces of the node}

\begin{center}\small
\begin{tabular}{lll}
\toprule
\textbf{Name} & \textbf{Type} & \textbf{Role} \\
\midrule
\multicolumn{3}{l}{\emph{Published}}\\
\topic{/tello/image\_raw} & \code{sensor\_msgs/Image} (rgb8) & camera, stamped on arrival at the PC \\
\topic{/tello/camera\_info} & \code{sensor\_msgs/CameraInfo} & intrinsics, same stamp as the image \\
\topic{/tello/telemetry} & \code{std\_msgs/String} (JSON) & Tello state + driver fields, 10 Hz \\
\topic{/tello/rc\_sent} & \code{geometry\_msgs/Twist} (body) & command really sent, 20 Hz \\
\multicolumn{3}{l}{\emph{Subscribed}}\\
\topic{/tello/cmd\_vel\_safe} & \code{Twist} (body, m/s, rad/s) & velocity requests \\
\topic{/ekf/pose} & \code{PoseWithCovarianceStamped} & pose and its uncertainty \\
\topic{/ekf/velocity} & \code{TwistStamped} (map) & velocity for the hold \\
\topic{/ekf/drift} & \code{Vector3Stamped} (map) & learnt drift, cancelled by the hold \\
\topic{/aruco/nearest\_marker} & \code{std\_msgs/Float32} & distance of the closest marker in view \\
\multicolumn{3}{l}{\emph{Services} (\code{std\_srvs/Trigger})}\\
\topic{\textasciitilde/takeoff} & \code{Trigger} & takeoff, then climb \\
\topic{\textasciitilde/land} & \code{Trigger} & land (never dropped) \\
\topic{\textasciitilde/emergency} & \code{Trigger} & motors off; \textasciitilde\ = \topic{/tello\_driver} \\
\bottomrule
\end{tabular}
\end{center}

The flight zone is not a topic. It is a \emph{file path} given to the driver as the parameter
\code{flight\_zone\_file}; Section~\ref{sec:zone-ros} explains who else reads that file.

\section{driver.py: the node}
\label{sec:driverpy}

\subsection{Module-level function}

\func{load\_camera\_info}{(path, width, height, frame\_id) -> (CameraInfo, calibrated: bool)}
Builds the \code{CameraInfo} message published with every image. With a calibration file (YAML written by
\file{tools/calibrate\_camera.py}, normally \file{config/tello\_camera.yaml}) it copies the matrix $K$, the
distortion vector $D$, the rectification $R$ (identity if absent) and the projection $P$ (if absent,
$P=[K\mid 0]$ built from $K$). If the file's \code{image\_width}/\code{image\_height} differ from the stream size it
raises \code{ValueError}: a calibration for another resolution would give wrong poses silently.

Without a file it returns an ideal pinhole with the measured horizontal field of view of the video stream,
$55.6^\circ$ (DJI's $82.6^\circ$ is the photo mode):
$$
f = \frac{W/2}{\tan(55.6^\circ/2)}, \qquad
K = \begin{bmatrix} f & 0 & W/2 \\ 0 & f & H/2 \\ 0 & 0 & 1 \end{bmatrix}, \qquad D = 0 .
$$
For the $960\times720$ stream, $f \approx 910.4$ px. The second return value tells the caller which case it was,
for the log line.

\subsection{The class TelloDriver}

\func{TelloDriver.\_\_init\_\_}{(self)}
The constructor does everything in a fixed order; each step matters.

\begin{enumerate}[leftmargin=*,itemsep=2pt]
\item \textbf{Parameters.} The main ones are listed in Table~\ref{tab:driver-params}. The model parameters default
to the constants of \file{core/model.py} ($K_x=0.66$, $K_y=0.83$, $\tau=0.25$ s, $L_u=0.35$ s, $K_z=0.9$).
\item \textbf{Zone.} \code{self.zone = load\_zone(flight\_zone\_file)} (Section~\ref{sec:zonepy}); if a zone was
loaded, its bounding rectangle is logged.
\item \textbf{Safety filter.} A \code{CommandFilter} is built with the zone, the gains, the model and all the
guard parameters; its warnings go to the ROS logger.
\item \textbf{Live parameters.} A callback \code{on\_params} is registered (see below).
\item \textbf{Duplicate check.} The node spins for 2 s, then asks the graph who already publishes
\topic{/tello/image\_raw}. If anyone does (typically a forgotten fake driver), it logs a fatal error and exits:
two image sources make the pose jump between two drones. The check is done \emph{before} its own publisher is
created, so it never finds itself.
\item \textbf{Publishers, subscriptions, services} of the table above. The three EKF and marker subscriptions are
one-line lambdas that pass the numbers straight into the filter (\code{on\_ekf\_velocity}, \code{on\_ekf\_drift},
\code{on\_nearest\_marker}).
\item \textbf{Drone object.} In mock mode a \code{FakeTello(mock\_markers\_file)}; if a zone exists and the
environment variable \code{FAKE\_START} is not set, the fake drone is moved to a start point inside the zone:
centre of the zone rectangle, pushed along $\bm n$, the horizontal direction pointing away from the main marker
wall, by $d=\max(0,\ |n_x|h_x+|n_y|h_y-0.3)$, where $h_x,h_y$ are the half-sizes of the rectangle. So the fake
drone starts 0.3 m from the zone edge furthest from the wall, facing the wall. In real mode
\code{wait\_for\_tello\_network} is called, the \code{djitellopy} logger is set to \code{WARNING} (otherwise it prints
\code{rc 0 0 0 0} twenty times per second) and a \code{Tello()} is created.
\item \textbf{Connect.} \code{tello.connect()}; on failure a fatal message explains the usual cause (another
program talked to the Tello first) and the process exits. Then the battery is logged and \code{streamon} is sent.
\item \textbf{Video reader.} \code{low\_latency\_video} (default \code{false}, forced off in mock mode) selects
\code{LowLatencyReader} (Section~\ref{sec:video}); otherwise \code{tello.get\_frame\_read()}, the \code{djitellopy}
reader (or the fake one).
\item \textbf{State and timers.} \code{rc\_target}=(0,0,0,0), \code{last\_cmd}=0, the lock \code{busy}, the
telemetry dict \code{state}, then the four timers of Table~\ref{tab:timers}.
\end{enumerate}

\begin{table}[ht]
\centering\small
\begin{tabular}{llp{5.8cm}}
\toprule
\textbf{Parameter} & \textbf{Default} & \textbf{Meaning} \\
\midrule
\code{mock} & false & fake drone (\file{sim/fake\_tello.py}) \\
\code{lin\_gain} & $0.38/30\approx0.01267$ & m/s per rc point (measured Nov 7) \\
\code{yaw\_gain} & $18^\circ/30\approx0.01047$ & rad/s per rc point \\
\code{max\_rc} & 80 & saturation of every rc axis ($\approx1.01$ m/s) \\
\code{cmd\_timeout} & 0.5 s & request older than this $\Rightarrow$ zero \\
\code{low\_battery} & 12 \% & land automatically at or below \\
\code{takeoff\_height} & 0 (launch: auto) & climb target after takeoff (m) \\
\code{max\_height} & 2.6 m & ceiling (room ceiling 2.80 m) \\
\code{flight\_zone\_file} & "" (launch: auto) & zone YAML; empty = no zone guards \\
\code{min\_wall\_distance}, \code{wall\_slow\_band} & 0.7 m, 0.15 m & wall guard \\
\code{blind\_speed} & 0.15 m/s & horizontal speed limit without a pose \\
\code{position\_hold}, \code{height\_hold} & true, true & hold on/off \\
\code{hold\_kp}, \code{hold\_kv}, \code{hold\_ki} & 0.8, 0.4, 0.15 & hold gains \\
\code{hold\_vmax}, \code{hold\_kp\_z}, \code{hold\_drift\_ff} & 0.4 m/s, 1.0, true & hold limits, drift cancel \\
\code{pose\_topic} & \topic{/ekf/pose} & pose for the hold and the guards \\
\code{low\_latency\_video} & false & use \file{core/video\_reader.py} \\
\code{video\_timeout}, \code{network\_timeout} & 3 s, 30 s & watchdog, Wi-Fi wait \\
\bottomrule
\end{tabular}
\caption{Parameters of \code{tello\_driver} (\file{core/driver.py}, lines 66--96 and 125--185).}
\label{tab:driver-params}
\end{table}

\begin{table}[ht]
\centering\small
\begin{tabular}{lll}
\toprule
\textbf{Timer} & \textbf{Period} & \textbf{Job} \\
\midrule
\code{poll\_video} & 10 ms (100 Hz) & publish a new image when the reader has one (30 fps max) \\
\code{poll\_state} & 100 ms (10 Hz) & telemetry JSON, low-battery landing \\
\code{send\_rc} & 50 ms (20 Hz) & the command loop: timeout, climb, safety chain, send \\
\code{video\_watchdog} & 1 s & restart the stream after 3 s without images \\
\bottomrule
\end{tabular}
\caption{The four timers of the driver.}
\label{tab:timers}
\end{table}

\begin{intuition}
\code{rclpy.spin} uses a single-threaded executor: the timers, subscriptions and service callbacks of the driver
never run at the same time. The only other threads are the ones of \code{djitellopy} (state and video receivers),
the video reader, and the short-lived threads that run blocking SDK commands (\code{run\_async}). The lock
\code{busy} is what keeps the 20 Hz rc loop quiet while such a command runs.
\end{intuition}

\func{on\_params}{(params) -> SetParametersResult}
A closure registered with \code{add\_on\_set\_parameters\_callback}. For every changed parameter it calls
\code{safety.set\_param(name, value)}; the names in \code{CommandFilter.LIVE\_PARAMS} (hold gains, hold on/off) are
applied immediately and logged, e.g. \code{ros2 param set /tello\_driver hold\_kp 1.0} during a flight. It always
returns success, so other parameters are accepted but have no effect after start-up.

\func{wait\_for\_tello\_network}{(timeout) -> None}
Waits until the PC has a route to the Tello. It "connects" a UDP socket to \code{192.168.10.1:8889}: for UDP this
sends nothing, it only makes the kernel pick a source address. If that address starts with \code{192.168.10.},
the PC is on the Tello's Wi-Fi. Otherwise it warns once and retries every second, and exits after
\code{network\_timeout} (30 s). The reason is WSL mirrored networking, which follows a Windows Wi-Fi change only
after a few seconds.

\func{video\_watchdog}{() -> None}
Runs at 1 Hz and acts only if no image arrived for \code{video\_timeout} = 3 s (never in mock mode, never while a
restart is running).
\begin{itemize}[itemsep=1pt]
\item Low-latency reader: that reader reopens the stream itself, so only \code{streamon} is re-sent.
\item \code{djitellopy} reader: it blocks on UDP port 11111 with no timeout. If its thread is alive, it will resume
by itself when packets come back, so only \code{streamon} is re-sent. If the thread is dead, the old container is
closed (which frees port 11111; a second reader on a busy port fails with "Address already in use"), and a new
reader is made after 0.5 s. This is done in a background thread so the executor is not blocked; the flag
\code{restarting} prevents overlapping restarts, and \code{last\_frame\_t} is reset so the stream gets 3 more seconds.
\end{itemize}

\func{poll\_video}{() -> None}
At 100 Hz, reads \code{reader.frame}. The readers return the \emph{same array object} until a newer image exists,
so \code{frame is self.last\_frame} detects "nothing new" without comparing pixels. The black $400\times300$
placeholder that \code{djitellopy} returns before the first decoded image is skipped. For a new image it updates
the frame-rate estimate (over 2 s windows, sent in telemetry as \code{video\_fps}), (re)builds the
\code{CameraInfo} when the image size changes, and publishes \code{Image} (\code{rgb8}, step $3W$) and
\code{CameraInfo} with the \emph{same} stamp: the current ROS time. The stamp is therefore the arrival time at
the PC, not the exposure time; the EKF accounts for the video latency $L_v=0.2$ s itself. In mock mode the
distortion is set to zero and passed to the fake camera (\code{world.set\_camera}), because the renderer draws an
ideal pinhole image.

\func{poll\_state}{() -> None}
At 10 Hz, copies \code{tello.get\_current\_state()} (the dict \code{djitellopy} parses from the 10 Hz telemetry
packets: \code{tof}, \code{h}, \code{bat}, \code{yaw}, \code{vgx}, \dots) into \code{self.state}, adds the driver's
own fields and publishes everything as JSON on \topic{/tello/telemetry}:
\begin{center}\small
\begin{tabular}{ll}
\toprule
\code{stamp} & ROS time in seconds \\
\code{flying} & \code{tello.is\_flying} \\
\code{busy} & a takeoff/land command is running \emph{or} the climb is not finished \\
\code{video\_fps} & images published per second \\
\code{hold} & horizontal position hold active \\
\code{guard} & the guard that changed the command in the last 0.5 s (\code{CommandFilter.guard\_now}) \\
\bottomrule
\end{tabular}
\end{center}
Missions read \code{busy} to wait for the end of the climb before sending velocities; \file{tools/flight\_report.py}
uses \code{hold} and \code{guard}. Finally, if the drone flies and the battery is at or below 12 \%, it starts a
landing with \code{run\_async} (repeated calls are harmless: \code{run\_async} does nothing while \code{busy} is held).

\func{on\_pose}{(m: PoseWithCovarianceStamped) -> None}
Converts the EKF pose to the numbers the filter wants: position, yaw from the quaternion,
$$
\psi = \operatorname{atan2}\big(2(q_wq_z+q_xq_y),\ 1-2(q_y^2+q_z^2)\big),
$$
and two standard deviations taken from the $6\times6$ covariance (row-major): $\sigma_{xy}=\sqrt{\max(\Sigma_{xx},\Sigma_{yy})}$
(indices 0 and 7) and $\sigma_z=\sqrt{\Sigma_{zz}}$ (index 14). Then \code{safety.on\_pose(x, y, z, yaw, sigma\_xy, sigma\_z)}.

\func{on\_cmd}{(m: Twist) -> None}
Converts a velocity request into Tello rc points and stores it with its reception time. The two conventions are
different:
\begin{center}\small
\begin{tabular}{lll}
\toprule
& \textbf{ROS body frame (REP 103)} & \textbf{Tello rc} \\
\midrule
order & $(v_x, v_y, v_z, \omega_z)$ & (left/right, forward/back, up/down, yaw) \\
positive & forward, \emph{left}, up, counter-clockwise & \emph{right}, forward, up, clockwise \\
\bottomrule
\end{tabular}
\end{center}
hence, with $g=$ \code{lin\_gain} and $g_\psi=$ \code{yaw\_gain} and $\operatorname{clip}$ = round and saturate to
$\pm$\code{max\_rc} $=\pm80$:
$$
\text{rc} = \Big(\operatorname{clip}\big(-\tfrac{v_y}{g}\big),\ \operatorname{clip}\big(\tfrac{v_x}{g}\big),\
\operatorname{clip}\big(\tfrac{v_z}{g}\big),\ \operatorname{clip}\big(-\tfrac{\omega_z}{g_\psi}\big)\Big).
$$
Example: $v_x=0.38$ m/s forward gives fb $=30$; $\omega_z=+0.314$ rad/s ($18^\circ$/s to the left) gives yaw $=-30$.
Nothing is sent here: the 20 Hz loop sends.

\func{send\_rc}{() -> None}
The heart of the node, at 20 Hz:
\begin{lstlisting}
if self.busy.locked():
    return                      # takeoff / land running: rc would interfere
rc = self.rc_target if time.time() - self.last_cmd < self.cmd_timeout else (0, 0, 0, 0)
flying = self.tello.is_flying
hgt, vz = self.safety.height_and_rate(self.state.get("tof"), flying)
if self.climb_target is not None and flying:
    ud = self.climb_rc(hgt, vz)
    if ud is not None:
        rc = (0, 0, ud, 0)      # climbing: other commands wait
rc = self.safety.filter(rc, hgt, vz, flying, busy=..., climbing=self.climb_target is not None)
self.tello.send_rc_control(*rc)
\end{lstlisting}
\begin{enumerate}[itemsep=1pt]
\item \textbf{Timeout.} A request older than 0.5 s is replaced by zero: if a mission crashes, the drone does not
keep its last velocity. (Zero is then turned into a position hold by the filter.)
\item \textbf{Height.} \code{height\_and\_rate} fuses the time-of-flight sensor (\code{tof}, cm) and the EKF
height, and differentiates it (Section~\ref{sec:safety}).
\item \textbf{Climb.} After a takeoff, while \code{climb\_target} is set, the request is replaced by a pure
vertical command from \code{climb\_rc}.
\item \textbf{Safety chain.} \code{safety.filter}.
\item \textbf{Send.} \code{send\_rc\_control} also keeps the link alive: the Tello lands by itself after 15 s
without any command.
\item \textbf{Publish what was sent} on \topic{/tello/rc\_sent}, converted back to the ROS body frame:
$v_x=g\,\text{fb}$, $v_y=-g\,\text{lr}$, $v_z=g\,\text{ud}$, $\omega_z=-g_\psi\,\text{yaw}$.
\end{enumerate}

\begin{keypoint}
The EKF's input is \topic{/tello/rc\_sent}, not \topic{/tello/cmd\_vel\_safe}. When a guard brakes, pushes back or
holds, the request and the real command differ; predicting with the request would make the EKF believe the drone
moves when it is braking.
\end{keypoint}

\func{stop\_rc}{() -> None}
Sends rc (0,0,0,0) once and publishes a zero \code{Twist} on \topic{/tello/rc\_sent}. The Tello keeps flying with
the last rc value while no new one comes, and none comes while a takeoff or landing runs (the 20 Hz loop is
paused by \code{busy}); flight 225048 flew on backwards for 1.2 m before landing. So every blocking command is
preceded by \code{stop\_rc}.

\func{run\_async}{(fn, label) -> None}
Runs a blocking SDK command (\code{tello.takeoff}, \code{tello.land}, which wait for the drone's "ok", several
seconds) in a daemon thread, so the executor keeps publishing video and telemetry. The worker takes
\code{busy} with a \emph{non-blocking} \code{acquire}: if another command already runs, the new one is dropped.
It then calls \code{stop\_rc}, runs \code{fn}, logs any exception and releases the lock.

\func{srv\_takeoff}{(req, res) -> res}
Refuses if the last telemetry battery is below 25 \% (if no telemetry was received yet the default is 100, i.e.
accepted). Otherwise starts \code{takeoff\_and\_climb} through \code{run\_async} and answers immediately ("taking
off, then climbing to 1.85 m").

\func{takeoff\_and\_climb}{() -> None}
Calls the SDK \code{takeoff} (the Tello rises to about 0.8 m and answers). If \code{takeoff\_height} $>0$, sets
$$
z_{\text{climb}} = \min\big(\text{takeoff\_height},\ h_{\max}-0.6\big) = \min(1.85,\ 2.0) = 1.85\ \text{m (band mode)},
$$
and the time origin of the climb. The climb itself is \emph{not} an SDK command: \code{up} was refused right after
takeoff ("error Not joystick", Nov 9), and the barometer height \code{h} was wrong (20 cm at 80 cm). So the climb
is done with joystick commands by \code{send\_rc}, closed-loop on the fused height. The value of
\code{takeoff\_height} comes from the launch file (Section~\ref{sec:zone-ros}): the middle of the zone's height
range, 1.85 m for the band zone, $(1.55+2.30)/2\approx1.92$ m for the 3D zone.

\func{climb\_rc}{(hgt, vz) -> int | None}
The vertical rc during the climb, \code{None} when the climb is over. With a height measurement:
$$
e = z_{\text{climb}} - (h + 0.8\,\dot h), \qquad \text{ud} = \operatorname{int}\big(\operatorname{sat}_{[-25,\,40]}(100\,e)\big),
$$
(\code{int} truncates towards zero),
a proportional law on the height predicted 0.8 s ahead (the drone keeps moving for about a second after the
command changes), with 1 rc point per cm of error, limited to $+40$ ($\approx0.5$ m/s) up and $-25$ down.
End conditions:
\begin{itemize}[itemsep=1pt]
\item done: $|z_{\text{climb}}-h|<0.10$ m and $|\dot h|<0.15$ m/s continuously for 0.5 s;
\item timeout: 12 s after takeoff, whatever the height;
\item no height at all: rc 0 for up to 10 s, then the climb is abandoned.
\end{itemize}
During the climb, the filter is told \code{climbing=True}: the height hold, the band guard and the zone guard are
off (the drone starts below the band), but the horizontal position hold and the ceiling stay on.

\func{end\_climb}{(msg) -> None}
Clears \code{climb\_target}, logs the reason and returns \code{None} (so \code{climb\_rc} can \code{return
self.end\_climb(...)}).

\func{srv\_land}{(req, res) -> res}
Clears the pending request and the climb. If no blocking command is running, lands through \code{run\_async}.
If one is (typically a takeoff still waiting for the Tello's answer), the landing must not be dropped: it calls
\code{stop\_rc} and sends the raw \code{land} command without waiting for the answer
(\code{send\_command\_without\_return}).

\func{srv\_emergency}{(req, res) -> res}
Clears the climb and calls \code{tello.emergency()}: the motors stop at once and the drone falls. Only for real
emergencies (the cage then hits the floor).

\func{shutdown}{() -> None}
Called by \code{main} on Ctrl+C or external shutdown. If the drone flies: rc zero and \code{land}. Then (real
drone only) a fire-and-forget \code{streamoff}, the video reader is stopped, and \code{\_\_del\_\_} of the
\code{Tello} class is replaced by a no-op: \code{djitellopy}'s destructor would otherwise call \code{end()}, which
retries \code{streamoff} for several seconds, even after the PC left the Tello's Wi-Fi.

\func{main}{() -> None}
\code{rclpy.init}, build the node, \code{spin} until Ctrl+C, then \code{shutdown} (landing), destroy and
\code{rclpy.shutdown}. The driver does not call \code{local\_ros\_env} itself: the launch file sets the same
environment variables (Section~\ref{sec:rosenv}).

\section{safety.py: the command filter}
\label{sec:safety}

\subsection{The chain}

\code{CommandFilter.filter} applies five stages, in this order, to the rc tuple (lr, fb, ud, yaw) that
\code{send\_rc} hands it:

\begin{center}\small
\begin{tabularx}{\linewidth}{l>{\raggedright\arraybackslash}X>{\raggedright\arraybackslash}p{2.2cm}}
\toprule
\textbf{Stage} & \textbf{What it does} & \textbf{Touches} \\
\midrule
\code{position\_hold} & no horizontal/yaw request: hold the point where the drone stops; no vertical request: hold the height & lr, fb, yaw; ud \\
\code{ceiling} & never above \code{max\_height} = 2.6 m, on the height predicted 1 s ahead & ud \\
\code{band\_guard} & not more than 10 cm out of the zone's height limits; pushed back 25 cm out & ud \\
\code{zone\_guard} & slow down / brake / push back before the obstacles; slow when blind & lr, fb \\
\code{wall\_guard} & a marker closer than 0.7 m ahead: no forward, back off & fb \\
\bottomrule
\end{tabularx}
\end{center}

The order matters. The hold comes first because it \emph{creates} commands (from zero) that the guards must still
check: a hold target that happens to be outside the zone is braked like any request. The guards that follow only
ever \emph{restrict} or push back; the last one to touch an axis wins. Each guard that acts calls \code{note(...)}
so the telemetry field \code{guard} tells which one changed the command.

The filter works in rc points at its input and output, but converts to metric velocities inside each stage with
the same \code{lin\_gain} $g$ and \code{yaw\_gain} $g_\psi$ as the driver. Three frames appear: the body frame
(forward $x_b$, left $y_b$), the map (world) frame, and rc. With yaw $\psi$ from the EKF,
$$
\begin{bmatrix} w_x \\ w_y \end{bmatrix} =
\begin{bmatrix} \cos\psi & -\sin\psi \\ \sin\psi & \cos\psi \end{bmatrix}
\begin{bmatrix} b_x \\ b_y \end{bmatrix},
\qquad b_x = g\,\text{fb},\quad b_y = -g\,\text{lr}.
$$

\subsection{The model used inside the filter}

All predictions use the identified horizontal model of \file{core/model.py} (drone with its cage, Nov 13 flights): per
axis, from a velocity command $u$ (m/s) to the real velocity $v$,
$$
\tau\,\dot v(t) = K\,u(t-L_u) - v(t), \qquad K_x = 0.66,\ K_y = 0.83,\ \tau = 0.25\ \text{s},\ L_u = 0.35\ \text{s},
$$
and $K_z = 0.9$ vertically. $K<1$ means the drone flies slower than the rc gain promises (the cage adds drag).
The zone guard uses the mean $\bar K = (K_x+K_y)/2 = 0.745$ because it works in the world frame where the direction
is arbitrary.

\subsection{Construction and inputs}

\func{CommandFilter.\_\_init\_\_}{(zone, gains, model\_*, guard and hold\_* parameters, warn, clock)}
Stores the parameters and initialises the state: no pose yet (\code{pose\_xy=None}), uncertainties 1.0 m, empty
histories, EKF velocity and drift zero, nearest marker "9 m" (far), no hold target, zero integrators. The
\code{clock} argument is injectable so an offline test can run faster than real time; everything in the filter
measures ages with it.

\func{on\_pose}{(x, y, z, yaw, sigma\_xy, sigma\_z) -> None}
Stores the pose and its time, and appends $(t,(x,y))$ to a 0.6 s history (for \code{velocity\_xy}). Before that,
it \emph{invalidates the hold targets} when the estimate cannot be trusted to mean the same point as before:
\begin{itemize}[itemsep=1pt]
\item horizontal target cleared if the new pose jumped more than 0.15 m from the previous one (markers back after
dead reckoning, EKF restart) or $\sigma_{xy}>0.06$ m;
\item height target cleared if $z$ jumped more than 0.15 m or $\sigma_z>0.08$ m.
\end{itemize}
Without this, a target taken from a drifting estimate would drag the drone to a meaningless point when the
markers come back (the docstring cites a 63 cm drag in a stress test).

\func{on\_ekf\_velocity}{(vx, vy, vz) -> None}
\func{on\_ekf\_drift}{(dx, dy) -> None}
\func{on\_nearest\_marker}{(d) -> None}
Store the EKF world velocity (with its time), the EKF drift estimate (the horizontal velocity the drone has with
zero command, learnt while markers are in view), and the distance of the closest marker in view (with its time).

\func{set\_param}{(name, value) -> bool}
If \code{name} is a key of \code{LIVE\_PARAMS} (\code{hold\_kp}, \code{hold\_kv}, \code{hold\_ki},
\code{hold\_kp\_z}, \code{hold\_vmax}, \code{hold\_drift\_ff}, \code{position\_hold}, \code{height\_hold}), sets the
mapped attribute and returns \code{True}.

\subsection{Helpers}

\func{pose\_valid}{(max\_age, max\_sigma) -> bool}
True if a pose exists, is younger than \code{max\_age} and has $\sigma_{xy}<$ \code{max\_sigma}. The EKF keeps
publishing without markers, with a growing covariance; this is where "dead reckoning for too long" becomes "no
pose". The thresholds differ per user: hold $(0.4\text{ s}, 0.12\text{ m})$, zone guard $(0.5\text{ s},
0.25\text{ m})$.

\func{velocity\_xy}{() -> (vx, vy)}
Finite difference between the oldest and newest pose of the 0.6 s history, $\big(\bm p_{\text{new}}-\bm
p_{\text{old}}\big)/\Delta t$, or zero if $\Delta t\le0.2$ s. Used by the hold when the EKF velocity is older than
0.3 s, and always by the zone guard.

\func{height\_estimate}{(tof\_cm) -> float | None}
The height above the floor used by the ceiling and the climb:
$$
h = \max\big(\text{tof}/100 \ \text{(if } 15\le\text{tof}\le400\text{ cm)},\ \ z_{\text{EKF}}\ \text{(if age} <0.5\text{ s and } \sigma_z<0.5\text{ m)}\big).
$$
Taking the \emph{larger} is conservative for the ceiling: over the bed the distance sensor reads about 0.85 m less
than the height above the floor, and the EKF height protects against that. \code{None} if neither is valid.

\func{height\_and\_rate}{(tof\_cm, flying) -> (h, dh/dt)}
\code{(None, 0)} when not flying. Otherwise $h$ from \code{height\_estimate}, appended to a 0.6 s history, and
$\dot h$ by finite difference over that history (zero if it spans less than 0.2 s).

\func{clip}{(v) -> int}
\code{round} then saturate to $[-80, 80]$.

\func{note}{(s) -> None}
\func{guard\_now}{() -> str}
\func{warn}{(key, msg, period=1.0) -> None}
\code{note} records which guard acted and when; \code{guard\_now} returns it if it is younger than 0.5 s (else
""); \code{warn} logs a message at most once per \code{period} per \code{key}, since the guards run 20 times per
second.

\subsection{Position hold}

\begin{intuition}
The Tello normally holds its position with a downward camera (optical flow). On the plain floor of the bedroom
that does not work: the drone reported zero velocity while drifting 0.4 m in 8 s. So when nobody asks for motion,
the driver holds the position itself, from the EKF pose. Three tricks make the hold stiff without oscillating:
(1) control the position the drone \emph{will} have when the command takes effect, (2) compute the command by
inverting the identified model instead of a plain PID, (3) subtract the drift that the EKF has learnt.
\end{intuition}

\func{hold\_axis}{(k, target, pos, vel, drift, K, kp, dt) -> u}
One axis of the hold. With $x^\ast$ the target, $x$, $v$ the position and velocity, $d$ the drift:
\begin{align*}
e &= x^\ast - (x + v\,L_u) && \text{error on the position predicted } L_u = 0.35\text{ s ahead}\\
v_d &= \operatorname{sat}_{\pm v_{\max}}(k_p\,e) && v_{\max}=0.4\text{ m/s},\ k_p = 0.8\ \text{s}^{-1}\\
i_k &\leftarrow \operatorname{sat}_{\pm0.1}(i_k + k_i\,e\,\Delta t)\quad\text{only if } |e|<0.3\text{ m} && k_i = 0.15\\
u &= \operatorname{sat}_{\pm0.6}\Big(\frac{v_d - d + k_v\,(v_d - v)}{K} + i_k\Big) && k_v = 0.4
\end{align*}
Reading the last line: in steady state the model says $v = K u + d$; asking for $v = v_d$ gives the
\emph{model inversion} $u=(v_d-d)/K$. The term $k_v(v_d-v)$ corrects the velocity error that remains (wrong $K$,
gusts), and the small integral $i_k$ (at most 0.1 m/s, built only near the target so a long transit does not wind
it up) removes the last static error. $\Delta t$ is the time since the previous call, at most 0.1 s. For $x,y$ the
function is called with $K=1$: the result is the \emph{world} velocity wanted, and the real inversion is done after
the rotation to the body frame, where the gains differ.

\func{position\_hold}{(rc, flying, busy=False, climbing=False) -> rc}
\begin{enumerate}[itemsep=1pt]
\item Off (both targets cleared, rc unchanged) if the hold is disabled, the drone is not flying, or
\code{busy}.
\item Any non-zero lr, fb or yaw: the pilot or the mission is moving, the horizontal target and integrators are
cleared and the request passes. Any non-zero ud, or the climb: same for the height.
\item No pose valid at $(0.4\text{ s}, 0.12\text{ m})$: rc unchanged (no hold, integrators kept).
\item Velocity: the EKF one if younger than 0.3 s, else \code{velocity\_xy} (and $v_z=0$). Drift: the EKF one if
\code{hold\_drift\_ff} and the EKF velocity is fresh, else zero.
\item \textbf{Horizontal hold}, if no horizontal request and (a target exists or $\sigma_{xy}<0.06$ m). A new target
is only taken on a certain pose, and not where the keys were released but where the drone \emph{will stop on its
own}. With zero command it keeps its speed for $L_u$, then decays with time constant $\tau$, so it travels
$$
\int_0^{L_u} v\,dt + \int_0^\infty v\,e^{-t/\tau}dt = (L_u+\tau)\,v = 0.6\,v ,
\qquad
(x^\ast,y^\ast,\psi^\ast) = (x+0.6\,v_x,\ y+0.6\,v_y,\ \psi).
$$
Then, per world axis, $w_x=$ \code{hold\_axis}$(x)$, $w_y=$ \code{hold\_axis}$(y)$; rotation into the body frame and
inversion of the per-axis gain:
$$
b_x = \operatorname{sat}_{\pm0.6}\frac{\cos\psi\,w_x+\sin\psi\,w_y}{K_x},\qquad
b_y = \operatorname{sat}_{\pm0.6}\frac{-\sin\psi\,w_x+\cos\psi\,w_y}{K_y},
$$
$$
\omega =\operatorname{sat}_{\pm0.3}\big(0.8\,\operatorname{wrap}(\psi^\ast-\psi)\big),
$$
and back to rc: lr $=$ clip$(-b_y/g)$, fb $=$ clip$(b_x/g)$, yaw $=$ clip$(-\omega/g_\psi)$. A body command of
0.6 m/s is about 47 rc points.
\item \textbf{Height hold}, if enabled and no vertical request and not climbing. New target
$z^\ast = z + 0.5\,v_z$ (where it stops), then $u_z=$ \code{hold\_axis}$(2, z^\ast, z, v_z, 0, K_z{=}0.9, k_{pz}{=}1.0)$,
limited to $\pm0.3$ m/s, i.e.\ $\pm24$ rc points.
\end{enumerate}
The docstring reports the tuning on the fake drone with a drift of 8/5 cm/s: the hold stays within 2.4 cm and
overshoots 7--9 cm after a key release (the previous PI: 13 cm and 24 cm).

\subsection{Ceiling}

\func{ceiling}{(rc, hgt, vz) -> rc}
Does nothing without a height. Otherwise, with $h_{\max}=2.6$ m, the limit $\ell = h_{\max}-0.15 = 2.45$ m and the
height predicted 1 s ahead (only when rising) $\hat h = h + \max(\dot h,0)\cdot1\,\text{s}$:
$$
\text{ud} \leftarrow
\begin{cases}
-30 & h > h_{\max} \quad\text{(descend, whatever was asked)}\\
\min(\text{ud}, -20) & \hat h > \ell \quad\text{(brake now)}\\
\min\big(\text{ud},\ \operatorname{int}((\ell-\hat h)/g)\big) & \text{otherwise.}
\end{cases}
$$
The last line limits the commanded climb rate to $(\ell-\hat h)/(1\text{ s})$: the drone slows down
progressively as it approaches 2.45 m. One second is used because the drone needs about a second to stop
vertically. The ceiling acts also during the climb.

\subsection{Height band (band\_guard)}

\func{band\_guard}{(rc, flying, climbing=False) -> rc}
The obstacle grid of the zone is computed for a range of heights only (the band file says: 0.25 m from any
surface between 1.50 and 2.20 m), so the horizontal guard is only valid if the drone stays near that range. This
guard keeps it there. It acts only with a zone, in flight, not during the climb, and only on a fresh ($<0.5$ s),
certain ($\sigma_z\le0.10$ m) EKF height; markers are not required (an EKF height stays good for a while on the
distance sensor).

The limits $[z_{lo}, z_{hi}]$ come from \code{height\_limits(zone, x, y)}: the band $[z_{\min}, z_{\max}]$ of a 2D
zone, or, for the 3D zone, the floor and top of the boxes above the current point. The floor is "tight" when it
is higher than the zone's global floor (over the east furniture, $z_{lo}=2.08$ m, furniture top 1.88 m: only 0.20 m
of slack). With $h=z_{\text{EKF}}$:

\begin{center}\small
\begin{tabular}{lll}
\toprule
\textbf{Condition} & \textbf{Action} & \textbf{Note} \\
\midrule
$h < z_{lo} - 0.25$ (tight: $-0.12$) & ud $\ge 12$ & push up ($\approx0.15$ m/s) \\
$h < z_{lo} - 0.10$ (tight: $-0.03$) and ud $<0$ & ud $=0$ & no further descent \\
$h > z_{hi} + 0.25$ & ud $\le -12$ & push down \\
$h > z_{hi} + 0.10$ and ud $>0$ & ud $=0$ & no further climb \\
\bottomrule
\end{tabular}
\end{center}
In band mode ($1.70$--$2.00$ m) the pilot can therefore fly between 1.60 and 2.10 m, inside the 1.50--2.20 m for
which the grid is valid.

\subsection{Flight zone (zone\_guard, grid\_guard, box\_guard)}

\func{zone\_guard}{(rc, flying, climbing=False) -> rc}
Off without zone, on the ground or during the climb. Then:
\begin{itemize}[itemsep=1pt]
\item \textbf{Blind} (no pose valid at $(0.5\text{ s}, 0.25\text{ m})$): the zone cannot be enforced. Blocking all
motion was tried and trapped the pilot (he could not move to find the markers again); now every horizontal axis is
limited to $\text{lim}=\operatorname{round}(0.15/g)=12$ rc points, 0.15 m/s commanded. At that speed the pilot has
about a second to react over 15 cm. Yaw and vertical are not limited (turning is how markers are found).
\item \textbf{With a pose}: the horizontal request is converted to a world velocity $\bm w=(w_x,w_y)$ (body $\to$
world rotation above), passed to \code{grid\_guard} if the zone has a clearance grid, else to \code{box\_guard},
rotated back to the body frame and to rc. ud and yaw pass unchanged. The velocity used for predictions is
\code{velocity\_xy()}, the finite difference of poses.
\end{itemize}

\func{grid\_guard}{(x, y, vx, vy, wx, wy) -> (wx, wy)}
The main horizontal guard. It reads the clearance $c$ (distance to the nearest obstacle, Section~\ref{sec:zonepy})
and the unit direction $\hat{\bm u}$ pointing \emph{away} from that obstacle, at two points, both on the layer of
the current height $z$:
\begin{enumerate}[itemsep=1pt]
\item the current position: $d,\ \hat{\bm u}_0$;
\item the position predicted $H=1$ s ahead if the command $\bm w$ is kept, with the identified model (velocity
$\bm v$ kept during $L_u$, then first-order response to $\bar K\bm w$):
$$
\hat{\bm p} = \bm p + \bm v L_u + \bar K\bm w\,(H-L_u) + (\bm v-\bar K\bm w)\,f,\qquad
f = \tau\big(1-e^{-(H-L_u)/\tau}\big) \approx 0.231\ \text{s}.
$$
This is the exact integral of $\tau\dot v = \bar K w - v$ over $[L_u, H]$ starting from $v$: a deceleration
already commanded is not taken for a run at the wall.
\end{enumerate}
The worse of the two is used: $(d_u, \hat{\bm u}) = (d_p, \hat{\bm u}_p)$ if $d_p<d$, else $(d,\hat{\bm u}_0)$. If
$\hat{\bm u}=\bm 0$ (flat clearance), nothing is done. The speed towards the obstacle is
$s = -\bm w\cdot\hat{\bm u}$ ($>0$: approaching). With the margin $m$ (0.25 m for the band and 3D zones):

\begin{center}\footnotesize
\begin{tabular}{lll}
\toprule
\textbf{Case} & \textbf{Correction} & \textbf{Result} \\
\midrule
$d < m$ (already too close) & if $s>-0.25$: $\bm w \mathrel{+}= (s+0.25)\,\hat{\bm u}$ & moves away at $\ge0.25$ m/s \\
$d_u < m+0.10$ (will be too close) & if $s>-0.15$: $\bm w \mathrel{+}= (s+0.15)\,\hat{\bm u}$ & moves away at $\ge0.15$ m/s \\
otherwise & if $s>a$, $a=(d_u-m)/1\text{ s}$: $\bm w \mathrel{+}= (s-a)\,\hat{\bm u}$ & approaches at $\le a$ \\
\bottomrule
\end{tabular}
\end{center}

\begin{intuition}
Only the component of the velocity along $\hat{\bm u}$ (towards or away from the nearest obstacle) is changed; the
tangential component is left alone, so the drone \emph{slides along} a wall instead of stopping dead. Far from the
obstacles the allowed approach speed is the remaining distance per second: with 0.5 m of room beyond the margin,
0.5 m/s towards the obstacle is allowed; it shrinks linearly to zero as the predicted point reaches the margin;
10 cm before it, the drone is braked into a retreat; inside the margin it is pushed out. Note that the speeds are
\emph{commanded} speeds, which the drone executes with the gain $K<1$.
\end{intuition}

\func{box\_guard}{(x, y, vx, vy, wx, wy) -> (wx, wy)}
The fallback when the zone file has no grid: the rectangle $[x_{\min},x_{\max}]\times[y_{\min},y_{\max}]$ only,
with a fixed margin $m_b=0.15$ m and a cruder prediction $\hat p = p + v\cdot 1\,$s. Per axis and per side, e.g.
the upper side:
$$
w \leftarrow
\begin{cases}
\min(w,-0.25) & p > \text{hi}\\
\min(w,-0.15) & \hat p > \text{hi}-m_b\\
\min\big(w,\ (\text{hi}-m_b-\hat p)/1\text{ s}\big) & \text{otherwise,}
\end{cases}
$$
symmetrically for the lower side with $\max$. With the shipped zone files this function is not used (all three
have a grid, explicitly or built from the boxes).

\subsection{Wall ahead}

\func{wall\_guard}{(rc, flying) -> rc}
A guard that needs neither the map nor the pose: the localizer publishes on \topic{/aruco/nearest\_marker} the
distance of the closest marker seen, from its apparent size, $d = f_x\,s/\ell_{\text{px}}$ ($s$ the printed side,
$\ell_{\text{px}}$ its longest edge in the image). The camera looks forward, so a big marker means a wall ahead.
Active if flying and the measurement is younger than 0.4 s:
$$
\text{fb} \leftarrow
\begin{cases}
\min(\text{fb}, -12) & d < 0.7\text{ m (back off at}\approx 0.15\text{ m/s)}\\
\operatorname{int}\Big(\text{fb}\,\dfrac{d-0.7}{0.15}\Big) & 0.7 \le d < 0.85\text{ m and fb}>0\\
\text{fb} & \text{otherwise.}
\end{cases}
$$
It is last in the chain so it has the final word on the forward axis. It is the backup for the case where the
pose is wrong but the camera is not.

\func{filter}{(rc, hgt, vz, flying, busy=False, climbing=False) -> rc}
\code{wall\_guard(zone\_guard(band\_guard(ceiling(position\_hold(rc)))))}, with the arguments each one needs.

\section{zone.py: the flight zone}
\label{sec:zonepy}

\subsection{What a clearance grid is, and what it is for}

The room contains walls, a bed, a desk and tall furniture along the east wall. The phone scan of the room
(\file{tools/room\_mesh\_to\_map.py}, aligned to the \topic{map} frame) gives all those surfaces as a triangle mesh.
The guard needs to answer, 20 times per second, one question: \emph{how far is the nearest obstacle from this
point, and in which direction should I move to get away from it?} A \textbf{clearance grid} answers it by a single
array lookup.

It is a 2D array over the floor plan, cell size $\delta$ (5 cm for the scan grids, 2 cm for the 3D zone), with
cell $(i,j)$ at $x = x_0 + j\delta$, $y = y_0 + i\delta$ (rows are $y$, columns are $x$). Each cell holds
$$
c(\bm p) = \min_{\bm q \in \mathcal O} \lVert \bm p - \bm q \rVert ,
$$
the Euclidean distance from the cell to the nearest obstacle cell $\mathcal O$ (a cell where some scan point lies
in the height range of the zone, or the outside of the room). Computing that for every cell at once is a
\textbf{Euclidean distance transform} (EDT): given a binary image, it returns for every pixel the distance to the
nearest pixel of the other class. OpenCV (\code{cv2.distanceTransform}, used by \file{tools/free\_zone.py}) and
SciPy (\code{scipy.ndimage.distance\_transform\_edt}, used by \file{core/zone.py}) both implement it in linear time.

The gradient $\nabla c$ points in the direction where the distance to obstacles grows fastest, i.e.\ \emph{away
from the nearest obstacle}; for an exact distance field $\lVert\nabla c\rVert=1$ almost everywhere, and the
direction is $\hat{\bm u} = \nabla c/\lVert\nabla c\rVert$. So the clearance is used for:
\begin{itemize}[itemsep=1pt]
\item the \textbf{driver's guard} (\code{grid\_guard}): how much room is left ($c - m$) and which direction is
"towards the obstacle" ($-\hat{\bm u}$);
\item the \textbf{missions' planning}: a waypoint or a takeoff point is acceptable only if
$c \ge m + $ a few cm (e.g.\ \file{missions/tour.py} uses $m+0.15$ m, \file{missions/patrol.py} $m+0.1$ m);
\item \textbf{RViz}: \topic{/room/free\_space}, an occupancy grid with 100 where $c<m$.
\end{itemize}
The \textbf{margin} $m$ is the clearance under which a point is no-fly: the drone with its cage is about 0.2 m
wide, plus pose error.

\begin{keypoint}
Instead of storing obstacles and computing distances on the fly, the zone precomputes the distance field once.
The guard then has the same cost whatever the shape of the room, and gets the escape direction for free from the
gradient. A rectangle check would forbid the whole L-shaped free space except its biggest rectangle; the clearance
grid lets the drone use any free shape.
\end{keypoint}

\subsection{Three zone files}

\begin{center}\small
\begin{tabular}{lllll}
\toprule
\textbf{File (config/)} & \textbf{Rectangle} $x\times y$ (m) & \textbf{Heights} & \textbf{Margin} & \textbf{Grid} \\
\midrule
\file{flight\_zone.yaml} & $[-0.26,0.79]\times[0.51,2.71]$ & 1.40--2.40 & 0.35 & \code{\_grid.npz} \\
\file{flight\_zone\_band.yaml} & $[-0.86,0.94]\times[0.36,2.76]$ & 1.70--2.00 & 0.25 & \code{\_band\_grid.npz} \\
\file{flight\_zone\_3d.yaml} & union of 2 boxes & 1.55--2.30 & 0.25 & from the boxes \\
\bottomrule
\end{tabular}
\end{center}

The first two are produced by \file{tools/free\_zone.py} from the scan: it rasterises every scan point between
the given heights (1.2--2.6 m for the first, 1.5--2.2 m for the band) into 5 cm cells, flood-fills the inside of the
room, runs the distance transform, finds the largest free rectangle (the "mission rectangle" \code{x\_min}\dots) and
removes narrow passages so that only the main free region keeps a clearance above the margin. The \code{.npz}
contains \code{clearance} ($99\times75$ cells, origin $(-1.56, -1.24)$, cell 0.05 m), \code{x0}, \code{y0},
\code{cell}, \code{z\_lo}, \code{z\_hi}. The rectangle and \code{z\_min}/\code{z\_max} in the YAML are the
\emph{allowed} band (1.70--2.00 for band mode), narrower than the heights the grid was computed for (1.50--2.20),
so the 10 cm tolerance of \code{band\_guard} stays inside the valid range.

The third file is drawn by hand (with \file{tools/zone\_view.py}) as a union of boxes:
\begin{lstlisting}
margin: 0.25
boxes:
  - {name: main,                 x: [-0.94, 0.67], y: [0.30, 2.87], z: [1.55, 2.30]}
  - {name: above_east_south,     x: [0.67, 1.15],  y: [0.30, 1.55], z: [2.08, 2.30]}
  - {name: above_east_north,     x: [0.67, 1.35],  y: [1.55, 2.87], z: [2.08, 2.30]}
\end{lstlisting}
Each box is where the drone \emph{centre} may be, already 0.25 m away from every scanned surface (walls at
$x=-1.24$ and $1.65$, $y=0$ and $3.17$; east furniture front at $x=0.92$, top 1.88 m). Seen from the south,
the union is an L: the whole room width between 2.08 and 2.30 m, but only the part west of $x=0.67$ down to
1.55 m. This lets the drone fly lower (more markers in view on the low rows) away from the furniture, and above
the furniture only high.

\subsection{Functions}

\func{wrap}{(a) -> float}
$\operatorname{atan2}(\sin a, \cos a)$: an angle brought back to $(-\pi,\pi]$.

\func{make\_grid}{(c, x0, y0, cell, margin) -> dict}
Turns a clearance array into the dict that \code{clearance()} reads:
\code{\{"d": c, "gx", "gy", "x0", "y0", "cell"\}}. The direction field is first the numerical gradient,
\code{gy, gx = np.gradient(c)} (axis 0 is $y$). Problem: inside an obstacle or outside the scan $c$ is a flat 0,
so the gradient is zero and the guard would not push at all (a stress test left the drone at $x=2.2$ m). Fix:
\begin{enumerate}[itemsep=1pt]
\item \code{bad} = cells with $c < m + 0.10$;
\item one EDT on \code{bad} with \code{return\_indices=True} gives, for each bad cell, the indices $(i^\ast, j^\ast)$
of the nearest good cell;
\item on the cells that are \code{bad} with a zero gradient, \emph{or} have $c<m$, the direction is replaced by
$(j^\ast-j,\ i^\ast-i)$: the vector to the nearest cell with room. It is not normalised here;
\code{clearance()} normalises.
\end{enumerate}
(In \code{flat = bad \& (...) | (c < margin)}, \code{\&} binds tighter than \code{|}.) So inside the no-fly part,
the push always points to the nearest flyable place, even where the gradient is flat or points into a dead end.

\func{box\_layers}{(boxes, margin, cell=0.02, pad=0.6) -> [(z\_lo, z\_hi, grid)]}
Converts the union of boxes into clearance grids, one per height layer.
\begin{enumerate}[itemsep=1pt]
\item A common 2D raster covering all boxes plus 0.6 m of padding, 2 cm cells ($176\times190$ cells for the 3D
file).
\item The \emph{layers} are the intervals between consecutive distinct box floors/tops: here
$\{1.55, 2.08, 2.30\}$ gives $[1.55, 2.08]$ and $[2.08, 2.30]$.
\item For each layer, \code{inside} = union of the footprints of the boxes that cover the \emph{whole} layer
height: the lower layer contains only \code{main}, the upper one \code{main} $\cup$ \code{above\_east\_furniture}.
\item The clearance is built so that the box boundary is the level $c=m$:
$$
c = \begin{cases}
m + \delta\cdot\mathrm{EDT}(\text{inside}) & \text{inside the union (distance to its edge)}\\
\max\big(0,\ m - \delta\cdot\mathrm{EDT}(\text{outside})\big) & \text{outside (distance to the union)}
\end{cases}
$$
Since each box is already $m$ from the scanned surfaces, $c$ approximates "distance to the nearest obstacle"
near the boundary, and the same guard code works unchanged.
\item \code{make\_grid} adds the directions.
\end{enumerate}
Evaluated on the shipped file: $c(0, 1.5)$ on the lower layer is 0.93 m with $\hat{\bm u}=(-1,0)$ (the nearest
edge is the east side of \code{main}, at $x=0.67$); at $(1.0, 1.5)$ the lower layer gives $c=0$,
$\hat{\bm u}=(-1,0)$ (push west, out of the furniture area), while the upper layer gives 0.61 m (free).

\begin{keypoint}
\code{box\_layers} runs only inside the Python processes that call \code{load\_zone}: the driver (enforcement) and
the missions (planning). RViz does not use it. \file{marker\_map\_publisher.py} reads the same YAML with
\code{yaml.safe\_load} and draws each box directly as a translucent cube with thick edges on
\topic{/room/flight\_zone}. The boxes you see in RViz and the boxes the driver enforces come from the \emph{same
file}, but through two different code paths.
\end{keypoint}

\func{load\_zone}{(path) -> dict | None}
Empty path: \code{None} (no zone, the zone and band guards are off). Otherwise two formats:
\begin{itemize}[itemsep=1pt]
\item \textbf{boxes} (3D): the rectangle keys are the bounding box of the boxes (for the file above
$x\in[-0.94,1.35]$ (v1.1: three boxes, the bounding box is unchanged), $y\in[0.30,2.87]$, $z\in[1.55,2.30]$), \code{margin} (default 0.25), the boxes as floats,
\code{layers} from \code{box\_layers}, and \code{grid} = the layer with the fewest cells at $c\ge m$, i.e.\ the
most restrictive one (the lower layer here), used when a caller gives no height.
\item \textbf{rectangle} (2D): the six keys, \code{margin} (default 0.35) and, if \code{grid\_file} is set, the grid
loaded from the \code{.npz} next to the YAML and passed through \code{make\_grid}.
\end{itemize}
Any exception (missing file, bad YAML, missing key) returns \code{None}. Note the consequence in the review.

\func{grid\_at}{(zone, z=None) -> grid}
The grid valid at height $z$. 2D zone or $z$ unknown: \code{zone["grid"]}. 3D zone: $z$ is clamped to
$[1.55, 2.30]$ and the highest layer whose floor is $\le z + 0.05$ is chosen. The 5 cm slack under a layer floor
avoids flipping between layers because of height noise around 2.08 m: from $z\ge2.03$ m the upper layer applies.

\func{height\_limits}{(zone, x, y) -> (lo, hi)}
For a 2D zone, $(z_{\min}, z_{\max})$. For a 3D zone, the boxes whose footprint contains $(x,y)$ grown by 5 cm;
the result is $(\min z_{lo}, \max z_{hi})$ over them. At $(0,1.5)$: $(1.55, 2.30)$; at $(1.0,1.5)$: $(2.08, 2.30)$;
on the shared edge $x=0.67$ both boxes count, so the lower floor wins. Outside every footprint (being pushed back
in), the limits of the nearest box(es) are used.

\func{clearance}{(zone, x, y, z=None) -> (d, (ux, uy)) | None}
The lookup used by everybody: picks the grid with \code{grid\_at}, converts $(x,y)$ to the nearest cell
$j=\operatorname{round}((x-x_0)/\delta)$, $i=\operatorname{round}((y-y_0)/\delta)$, returns $c_{ij}$ and the
normalised direction $(g_x, g_y)/\lVert(g_x,g_y)\rVert$ (zero vector if the norm is below $10^{-6}$). Outside the
array, it returns $d=0$ and the unit vector towards the centre of the zone rectangle. \code{None} if the zone has
no grid.

\section{How the zone exists in ROS}
\label{sec:zone-ros}

\emph{"We delimited a zone. How? Is it only in the driver node? How does it work in ROS?"}

The zone is three things: \textbf{data} (a YAML file, possibly with an \code{.npz}), \textbf{code} that interprets it
(\file{core/zone.py}, a plain module), and a \textbf{parameter} that tells each node which file to use. There is no zone
topic and no zone server. The path is chosen once, in \file{real.launch.py}:

\begin{lstlisting}
def zone_file(arg, band=False):
    if arg == "none":
        return ""
    if arg == "3d":   # the L-shaped zone (union of boxes)
        return str(Path(REPO) / "config" / "flight_zone_3d.yaml")
    if arg != "auto":
        return arg    # any path
    p = Path(REPO) / "config" / ("flight_zone_band.yaml" if band else "flight_zone.yaml")
    return str(p) if p.exists() else ""
\end{lstlisting}

With the defaults (\code{band:=true}, \code{flight\_zone:=auto}) the file is \file{config/flight\_zone\_band.yaml};
\code{flight\_zone:=3d} selects the L zone; \code{flight\_zone:=none} disables it. The same string is passed as the
parameter \code{flight\_zone\_file} to two nodes of the launch, and \code{takeoff\_height} is derived from it
(\code{takeoff\_height()}: in band mode the middle of the file's heights, or of the boxes' full range).

\begin{center}\small
\begin{tabularx}{\linewidth}{>{\raggedright\arraybackslash}p{3.8cm}>{\raggedright\arraybackslash}X>{\raggedright\arraybackslash}p{4.6cm}}
\toprule
\textbf{Who} & \textbf{How it gets the file} & \textbf{What it does with it} \\
\midrule
\code{tello\_driver} & parameter \code{flight\_zone\_file} from the launch & \code{load\_zone} $\to$ \code{CommandFilter}: \textbf{enforcement} (band and zone guards) \\
\code{marker\_map\_publisher} & same parameter from the launch & RViz only: boxes on \topic{/room/flight\_zone} (yellow: rectangle of a 2D file; green: the boxes of \file{flight\_zone\_3d.yaml}), \topic{/room/free\_space} from the \code{.npz} \\
missions (\code{Mission} and its subclasses: tour, goto, patrol, strike) & own parameter; if empty, \emph{asks the running driver} for its \code{flight\_zone\_file} (\code{GetParameters} service of the driver, function \code{driver\_zone\_file} in \file{missions/base.py}); else the default file & \code{load\_zone} again in their own process: \textbf{planning} (choosing targets with enough clearance, landing if the zone is left) \\
\code{tello\_teleop} & does not know the zone & protected by the driver \\
\bottomrule
\end{tabularx}
\end{center}

\begin{intuition}
Think of it as \emph{one file, read independently by several processes}. Enforcement is only in the driver; the
missions only use the zone to avoid \emph{asking} for something the driver would refuse (a mission that plans a
waypoint in the no-fly part would fight the guard forever). Asking the driver for its parameter guarantees both
read the same file, so a launch with \code{flight\_zone:=3d} switches enforcement, planning and display together.
\end{intuition}

\section{video\_reader.py: the low-latency reader}
\label{sec:video}

\subsection{Why}

\code{djitellopy}'s \code{BackgroundFrameRead} opens \code{udp://0.0.0.0:11111} with PyAV/FFmpeg defaults made for
smooth playback: input buffering, a long stream analysis at start, frame-threaded decoding (each thread holds
images back), no read timeout, and a huge UDP FIFO. When the decoder is slower than the 30 images/s the drone sends
(CPU busy with RViz and ArUco), packets pile up in that FIFO and the delay \emph{grows} without bound; on Nov 12 the
RViz image was a minute late, and a late image means a late pose for the EKF. The class
\code{LowLatencyReader} keeps the same interface (\code{.frame}, \code{.stop()}) but never lets the delay grow.

It is \textbf{off by default} (\code{low\_latency\_video:=false}): on Nov 13 it failed to open the real stream
(\code{InvalidDataError} in a loop); it was only validated on a synthetic H.264 stream (p90 latency 49 ms against
60 ms, max 68 against 120 ms).

\subsection{How}

\begin{itemize}[itemsep=1pt]
\item \textbf{URL} \code{udp://@0.0.0.0:11111?overrun\_nonfatal=1\&fifo\_size=8000}: the FIFO counts 188-byte units,
so 1.5 MB, about 1.5--3 s of 2--4 Mbit/s video; when full, data is dropped instead of delaying everything.
\item \textbf{Format options} \code{fflags=nobuffer}, \code{probesize=32768}, \code{analyzeduration=0}; codec flag
\code{low\_delay}; never frame threading (+260 ms measured).
\item \textbf{Timeouts} \code{(10 s open, 1 s read)}: after a Wi-Fi drop the read fails, the container is closed
(freeing the port) and reopened after 0.3 s by the reader's own thread.
\item \textbf{Only the newest image} is kept, and converted to RGB only when read.
\item \textbf{Catch-up}: the reader measures how long it waited for each packet.
\end{itemize}

\func{LowLatencyReader.\_\_init\_\_}{(url=URL, low\_latency=True, log=None, ...)}
Stores the settings (the defaults are the benchmarked ones; the other arguments exist for experiments), sets the
same black $400\times300$ placeholder as \code{djitellopy} (skipped by the driver), zero counters, and starts the
daemon thread \code{\_run}.

\func{frame}{(property) -> np.ndarray}
Takes the newest decoded \code{av.VideoFrame} under the lock (and clears it), converts it to an RGB array outside
the lock, and stores it as \code{\_frame} unless an even newer one was decoded meanwhile. Returns \code{\_frame}:
the \emph{same object} until a newer image exists, which is what \code{poll\_video}'s identity test relies on.

\func{stop}{(wait=2.0) -> None}
Raises \code{stopped} and joins the thread for up to 2 s. Only the reader thread may close the FFmpeg container
(closing it from another thread while it decodes crashes the process); it notices the flag within the 1 s read
timeout.

\func{stats}{() -> dict}
Frames decoded, mean decode time (ms), reconnections, skipped packets, number of catch-ups.

\func{\_close}{() -> None},\quad \func{\_open}{() -> (container, stream)},\quad \func{\_decode}{(packet) -> frames}
Close the container ignoring errors; open it with the options and set the codec flags; decode one packet (a method
so a test can override it).

\func{\_run}{() -> None}
The thread. In a loop until stopped: open; for each demuxed packet, record $(t_0, w)$ where $w$ is the time spent
waiting for that packet since the previous one was processed, over a sliding window of 1 s. The reader is declared
\textbf{behind} when the window holds more than 10 packets, spans more than 0.8 s, and
$$
\sum w < 0.03 \times 1\text{ s},
$$
i.e.\ it practically never had to wait for data: packets are queued and the images are late. It then
\textbf{skips} packets (counted in \code{skipped}) until the next key frame, or for at most 2 s (after which decoding
restarts with artefacts until the next key frame), and clears the window. A key frame is needed because H.264
P-frames only encode differences to the previous images. Decoded frames replace \code{\_av\_frame} under the lock.
Any exception (timeout, corrupt stream) is logged, the container closed, and the stream reopened after 0.3 s.

\section{ros\_env.py: local-only ROS}
\label{sec:rosenv}

\func{local\_ros\_env}{() -> None}
Must run before \code{rclpy.init}. Sets, only if not already set (\code{setdefault}), three variables:
\begin{itemize}[itemsep=0pt]
\item \code{ROS\_AUTOMATIC\_DISCOVERY\_RANGE=LOCALHOST},
\item \code{FASTDDS\_BUILTIN\_TRANSPORTS=LARGE\_DATA},
\item \code{FASTRTPS\_DEFAULT\_PROFILES\_FILE} = the installed \file{config/fastdds\_local.xml} (UDP on
127.0.0.1 only, packets $\le1400$ bytes).
\end{itemize} Why: all nodes run on one PC, and the PC switches its Wi-Fi to the Tello (no multicast,
no internet); WSL mirrored networking drops fragmented UDP on the loopback, and with Tailscale up the discovery
packets exceeded 1500 bytes and no node saw the others. Called by \code{tello\_teleop}, the missions
(\code{step\_test.run}) and \code{tello\_ekf}; the launch file sets the same three variables for the nodes it
starts (\code{LOCAL\_ROS\_ENV}), so a program started by hand in a second terminal must use the same settings to
see them.

\func{fastdds\_profile}{() -> str}
Path of \file{fastdds\_local.xml} in the package share directory, or "" if not found.

\section{The fake drone: fake\_tello.py and fake\_world.py}
\label{sec:fake}

With \code{mock:=true} the driver builds \code{FakeTello} instead of \code{djitellopy.Tello}. Everything downstream
(localizer, EKF, missions, RViz) is the real code; only the drone and its camera are simulated.

\subsection{fake\_tello.py}

\func{FakeFrameRead}{(world=None)}
The video reader. Its \code{frame} property produces a new image when the 30 Hz tick $k=\lfloor 30t\rfloor$ changes.
With a world, each tick renders the view (which also advances the dynamics, see \code{render}) and pushes it in a
queue; the image returned is the newest one at least $L_v=0.2$ s old, so the fake video has the same latency as the
real one (the EKF assumes it). Without a world: a grey image with a colour ramp and a moving white line.

\func{FakeTello}{(markers\_file="")}
The \code{djitellopy} subset the driver uses. \code{connect}, \code{streamon}, \code{streamoff}, \code{end}: no-ops;
battery 90 \%. \code{takeoff} sets \code{world.z\_target} $=\max(z, 0.85)$ and waits until reached (the real Tello
stops around 0.8--0.9 m); \code{land} targets 0.05 m and waits 1 s; \code{emergency} stops at once;
\code{send\_command\_without\_return("land")} lands in a thread; \code{send\_rc\_control} stores the rc in the
world; \code{get\_current\_state} returns the world's state (or constants without a world); \code{move\_up},
\code{get\_height}, \code{wait\_z} support the position commands.

\subsection{fake\_world.py}

\func{marker\_image}{(dictionary, mid, px=240) -> RGB image}
A printed marker: the ArUco square, $240$ px, centred on a white sheet 1.5 times larger.

\func{FakeWorld.\_\_init\_\_}{(markers\_file, width=960, height=720)}
Loads the real marker map (\file{config/markers.yaml}) with the localizer's own loader, renders one image per
marker, sets an ideal pinhole with the same $55.6^\circ$ field of view ($f\approx910$ px), the start pose, zero
velocity, an rc queue for the delay.

\func{set\_camera}{(K, D)}
Replaces the intrinsics with the driver's \code{CameraInfo} (distortion forced to zero by the driver), so the
renderer and the localizer use the same camera.

\func{start\_pose}{() -> (p, yaw)}
\code{FAKE\_START=x,y,yaw\_deg} if set. Otherwise: group the markers by the direction of their normal (rounded to
$45^\circ$), take the wall with the most markers (the west wall, 16 markers), stand on the floor 2 m in front of
its centre, facing it. The driver then moves this start inside the zone (Section~\ref{sec:driverpy}).

\func{step}{() -> None}
The dynamics, with $\Delta t \le 0.1$ s since the previous step:
\begin{itemize}[itemsep=1pt]
\item the rc applied is the one sent $L_u = 0.35$ s ago (queue);
\item body velocity target $\bm v^\ast = g\,(K_x\,\text{fb},\ -K_y\,\text{lr},\ \text{ud})$ and a first-order
Euler step $\bm v \mathrel{+}= (\bm v^\ast-\bm v)\min(1, \Delta t/\tau)$, $\tau=0.25$ s: the identified model;
\item during takeoff/land (\code{z\_target}), a position loop $v_z = \operatorname{sat}_{\pm0.6}(2\,\Delta z)$
that stops within 2 cm, like the Tello's own;
\item yaw rate first order towards $-g_\psi\,\text{yaw}$;
\item position integrated in the map with the yaw rotation; \code{FAKE\_DRIFT=vx,vy} adds a constant drift (the
real Tello without its optical-flow hold); $z\ge0.05$.
\end{itemize}

\func{render}{() -> RGB image}
Calls \code{step}, then draws: a vertical grey gradient (wall), then every marker, far ones first (painter's
algorithm). A marker is skipped if seen from behind or edge-on ($(\bm c-\bm t)\cdot\bm n \le 0.05$), or if a corner
is less than 0.15 m in front of the camera. Its four sheet corners are projected with \code{cv2.projectPoints}
(camera 5 cm ahead of the body, optical axes $R_{bc}$: $z$ forward, $x$ right, $y$ down), and the marker image is
warped onto that quadrilateral with a homography (\code{getPerspectiveTransform} + \code{warpPerspective}).
\code{FAKE\_BLACKOUT=a:b} returns a blank wall between $a$ and $b$ seconds after takeoff, to test the search
behaviour.

\func{state}{() -> dict}
A telemetry dict in Tello units and signs: \code{yaw} in degrees clockwise, \code{vgx}, \code{vgy}, \code{vgz} in
dm/s (body frame, $z$ down), \code{tof} and \code{h} in cm, battery 90, plus \code{true\_x}, \code{true\_y},
\code{true\_z}, \code{true\_yaw\_deg}: the ground truth, recorded in the session so the EKF can be scored.

\begin{intuition}
Because the dynamics advance only inside \code{render}, the simulation runs at the rate the driver polls the video
(30 steps per second). The fake drone obeys the same identified model the EKF and the hold assume, so the mock is
good for testing logic (guards, missions, search) but cannot reveal a model error; that is what the validation
flights are for.
\end{intuition}

\section{Code review notes}
\label{sec:driver-review}

\begin{review}
\begin{enumerate}[leftmargin=*,itemsep=3pt]
\item \textcolor{OliveGreen}{\textbf{[Fixed in v1.1, see Appendix~\ref{app:review}.]}} \textbf{A zone file that fails to load disables the zone silently.} \file{core/zone.py}:87--88 returns
\code{None} on \emph{any} exception, and \file{core/driver.py}:82--84 then simply logs nothing; the driver flies with
no band and no zone guard. Same for \file{real.launch.py}:77--78 (\code{auto} with a missing file gives ""). A
typo in a path or a corrupt \code{.npz} removes the main safety layer without a warning. Suggest: log an error and
refuse to start when \code{flight\_zone\_file} is non-empty but \code{load\_zone} fails.

\item \textbf{The horizontal hold's 0.12 m tolerance is dead.} \file{core/safety.py}:63--64 clears \code{hold\_target}
whenever $\sigma_{xy}>0.06$; \file{core/safety.py}:185 then holds only if a target exists or $\sigma_{xy}<0.06$. So
between 0.06 and 0.12 m (allowed by \code{pose\_valid(0.4, 0.12)}, line 176) the hold is never active. Either the
intent was to keep an existing target up to 0.12 m (then do not clear on $\sigma$ in \code{on\_pose}), or 0.12
should read 0.06.

\item \textbf{The height hold ignores $\sigma_z$.} \file{core/safety.py}:65--66 clears \code{hold\_zt} when
$\sigma_z>0.08$, but lines 201--205 immediately recreate it from the uncertain $z$ on the next cycle (the check
at line 176 is on $\sigma_{xy}$ only). The hold then follows the estimate instead of holding a height. Suggest a
$\sigma_z$ condition before taking a new height target.

\item \textbf{The zone guard uses the finite-difference velocity.} \file{core/safety.py}:275 calls
\code{velocity\_xy()} (0.6 s window, lags about 0.3 s, and a pose jump produces a large spurious velocity, hence a
spurious brake), while the hold uses the fresher EKF velocity when available (line 180). Using the EKF velocity
here too would make the 1 s prediction consistent with the hold.

\item \textbf{band\_guard trusts $(x,y)$ without checking $\sigma_{xy}$.} \file{core/safety.py}:235--242 checks the age
and $\sigma_z$, then calls \code{height\_limits(z, *self.pose\_xy)}. With the 3D zone, a wrong horizontal estimate
picks the wrong floor (1.55 m instead of 2.08 m over the furniture, or the reverse, pushing the drone up). Suggest
falling back to the global $z_{\min}$ of the most restrictive case (2.08 m near the boundary) when
$\sigma_{xy}$ is large.

\item \textbf{Mock start point is wrong for the 3D zone.} \file{core/driver.py}:141--146 uses the bounding rectangle.
With the main marker wall on the west, $\bm n=(1,0)$, and for \file{flight\_zone\_3d.yaml} the fake drone starts at
$x = 0.205+(1.145-0.3) \approx 1.05$ m, i.e.\ under the \code{above\_east\_furniture} box, where the lower layer is
no-fly: after the climb to 1.92 m the band guard pushes up and the zone guard pushes west at the same time. Use a
point with clearance (e.g.\ the maximum of \code{zone["grid"]}).

\item \textcolor{OliveGreen}{\textbf{[Fixed in v1.1, see Appendix~\ref{app:review}.]}} \textbf{A NaN request kills the node.} \file{core/driver.py}:325--326 calls \code{clip}, i.e.\ \code{round(v)}
(\file{core/safety.py}:125); \code{round(float("nan"))} raises \code{ValueError} (checked) inside the subscription
callback, which ends \code{rclpy.spin} (the drone then lands through \code{shutdown}). Reject non-finite values
(\code{math.isfinite}) in \code{on\_cmd}.

\item \textbf{Dead \code{busy} argument.} \file{core/driver.py}:332 returns while \code{busy} is locked, so the
\code{busy=self.busy.locked()} passed at line 341 is always \code{False}; the \code{busy} branch of
\code{position\_hold} (\file{core/safety.py}:165) is unreachable from the driver.

\item \textbf{Stale comments about the 3D zone.} \file{config/flight\_zone\_3d.yaml}:1--2 says "Not used by the driver
yet", but \file{real.launch.py}:73--74 enforces it with \code{flight\_zone:=3d} (commit \code{aab4639}); the
comment in \file{marker\_map\_publisher.py}:68 still calls it "proposed".

\item \textbf{Launch help disagrees with the code.} \file{real.launch.py}:158 says "no climbing above max-0.2"; the
ceiling uses $h_{\max}-0.15$ (\file{core/safety.py}:215).

\item \textbf{Obsolete re-exports.} \file{core/driver.py}:30 and 33 import \code{FakeFrameRead}, \code{clearance},
\code{wrap} "re-exported: missions import them here"; no module imports them from \code{driver} any more (the
missions import \code{tello\_driver.zone} directly). They can be removed.

\item \textbf{Fake vertical and yaw dynamics differ from \file{core/model.py}.} \file{sim/fake\_world.py}:195 uses a vertical
gain 1 (\code{ud * LIN\_GAIN}) while the hold inverts $K_z=0.9$; lines 191 and 205 apply $L_u=0.35$ s and
$\tau=0.25$ s to yaw, whereas the identified yaw model is $L=0.2$ s, $\tau\approx0$.

\item \textbf{The offline tests are not in the repository.} \file{docs/ARCHITECTURE.md}:62--63 and the docstrings of
\file{core/safety.py}/\file{core/zone.py} cite 400 randomized flights and numbered seeds, but no test file is tracked
(\code{git ls-files} lists none). The ROS-free design of \code{CommandFilter} is only worth it if that harness is
kept; commit it (e.g.\ \file{src/tello\_driver/test/}).

\item \textbf{No error handling in the emergency service.} \file{core/driver.py}:437--441 calls
\code{tello.emergency()} without \code{try}: if the link is down the service call fails with an exception instead of
an answer. Minor, but this is the one command that must always answer.
\end{enumerate}
\end{review}
